\documentclass[a4paper,12pt]{article}

\usepackage[T2A]{fontenc}
\usepackage[utf8]{inputenc}
\usepackage[russian]{babel}
\usepackage{amsmath,amssymb,amsthm,mathtools}
\usepackage{helvet}
\renewcommand{\familydefault}{\sfdefault}
\usepackage{icomma}
\usepackage[left=18mm,right=18mm,top=18mm,bottom=20mm]{geometry}
\usepackage{microtype}
\usepackage{tikz}
\usepackage{tkz-euclide}
\usepackage{pgfplots}
\pgfplotsset{compat=1.18}
\usepackage{tabularx,booktabs,longtable}
\usepackage{enumitem}
\usepackage[hidelinks]{hyperref}
\usepackage{parskip}
\usepackage{fancyhdr}
\usepackage{setspace}
\usepackage{xcolor}

\definecolor{accent}{HTML}{008080}
\definecolor{darkaccent}{HTML}{004D4D}
\definecolor{textgray}{HTML}{333333}
\definecolor{softgray}{HTML}{F3F6F6}

\providecommand{\tg}{\operatorname{tg}}
\providecommand{\ctg}{\operatorname{ctg}}
\providecommand{\arctg}{\operatorname{arctg}}

\newcommand{\Z}{\mathbb{Z}}
\newcommand{\R}{\mathbb{R}}
\newcommand{\important}[1]{\textcolor{darkaccent}{\textbf{#1}}}

\pagestyle{fancy}
\fancyhf{}
\lhead{\textcolor{textgray}{Логарифмические неравенства}}
\rhead{\textcolor{textgray}{ЕГЭ}}
\cfoot{\thepage}
\renewcommand{\headrulewidth}{0.4pt}
\renewcommand{\headrule}{\hbox to\headwidth{\color{accent}\leaders\hrule height \headrulewidth\hfill}}

\onehalfspacing
\setlist[itemize]{leftmargin=1.4em,itemsep=0.25em}
\setlist[enumerate]{leftmargin=1.6em,itemsep=0.35em}

\begin{document}

\begin{titlepage}
\thispagestyle{empty}
\centering
\vspace*{0.18\textheight}

{\Huge\bfseries\textcolor{darkaccent}{Чек-лист}\par}
\vspace{0.8em}
{\LARGE Логарифмические неравенства: ОДЗ, модули, замена\par}

\vspace{2.5em}
{\large Лёвин Артём Александрович эксклюзивно для Николь\par}

\vspace{1.2em}
{\large \today\par}

\vfill

\begin{tikzpicture}[x=1cm,y=1cm]
\draw[accent,line width=1.2pt]
  (-5.5,0) .. controls (-3.2,1.0) and (-1.1,-0.8) .. (1.2,0.15)
  .. controls (2.8,0.8) and (4.2,-0.35) .. (5.5,0.25);
\draw[darkaccent,line width=0.4pt]
  (-5.5,-0.25) .. controls (-2.8,0.55) and (1.8,-0.9) .. (5.5,-0.1);
\end{tikzpicture}
\end{titlepage}

\section*{1. Главная идея темы}

Логарифмические неравенства почти всегда решаются не одной формулой, а цепочкой аккуратных действий:
\[
\text{ОДЗ} \quad \longrightarrow \quad
\text{приведение логарифмов} \quad \longrightarrow \quad
\text{равносильный переход} \quad \longrightarrow \quad
\text{метод интервалов / замена}.
\]

Ключевой принцип: \important{нельзя механически «отбрасывать» логарифмы}, пока не учтены основание, подлогарифмические выражения и знак монотонности.

\section*{2. ОДЗ для логарифмов}

Для выражения
\[
\log_a f(x)
\]
обязательно:
\[
a>0,\qquad a\ne1,\qquad f(x)>0.
\]

Если основание является числом, например \(2\), \(4\), \(\frac14\), то условия на основание часто выполняются автоматически. Но условие \(f(x)>0\) проверяется всегда.

\subsection*{Автоматически выполняющиеся ограничения}

Иногда одно ограничение следует из другого условия. Например, если после перехода получено
\[
3-x\ge \frac{|x-2|}{x+1}
\]
и уже известно, что
\[
\frac{|x-2|}{x+1}>0,
\]
то \(3-x>0\) выполняется автоматически. В таком случае лишнее ограничение можно не решать отдельно.

\section*{3. Основание логарифма и смена знака}

Если \(a>1\), то логарифмическая функция возрастает:
\[
\log_a f(x)\ge \log_a g(x)
\quad \Longleftrightarrow \quad
f(x)\ge g(x),
\]
при \(f(x)>0,\ g(x)>0\).

Если \(0<a<1\), то логарифмическая функция убывает:
\[
\log_a f(x)\ge \log_a g(x)
\quad \Longleftrightarrow \quad
f(x)\le g(x).
\]

\important{При основании меньше единицы знак неравенства меняется.}

\section*{4. Модуль внутри неравенства}

\subsection*{Базовые переходы}

Если \(G(x)>0\), то:
\[
|F(x)|>G(x)
\quad \Longleftrightarrow \quad
\left[
\begin{array}{l}
F(x)>G(x),\\
F(x)<-G(x).
\end{array}
\right.
\]
Это \important{совокупность}, поэтому после решения двух неравенств берётся объединение.

Если \(G(x)\ge0\), то:
\[
|F(x)|\le G(x)
\quad \Longleftrightarrow \quad
-G(x)\le F(x)\le G(x).
\]
Это система, поэтому берётся пересечение условий.

\subsection*{Раскрытие модуля по областям}

Если модуль не раскрывается одинаково на всей ОДЗ, нужно рассмотреть случаи:
\[
|x-2|=
\begin{cases}
2-x, & x<2,\\
x-2, & x\ge2.
\end{cases}
\]

Если по ОДЗ точка \(x=2\) исключена, это учитывается в итоговом пересечении.

\section*{5. Метод интервалов}

После преобразований часто получается рациональное неравенство. Алгоритм:

\begin{enumerate}
\item Перенести всё в одну сторону.
\item Привести к общему знаменателю.
\item Разложить числитель и знаменатель на множители.
\item Отметить на оси нули числителя и нули знаменателя.
\item Нули знаменателя всегда выколоты.
\item Расставить знаки на промежутках.
\item Выбрать промежутки по знаку неравенства и пересечь с ОДЗ.
\end{enumerate}

\begin{center}
\begin{tikzpicture}[x=1.35cm,y=1cm]
\draw[->,line width=0.8pt] (-3.2,0) -- (3.4,0) node[right] {$x$};
\foreach \x/\lab/\type in {-2/$-2$/fill,0/$0$/draw,2/$2$/fill}{
  \ifx\type\fill
    \filldraw[darkaccent] (\x,0) circle (2pt);
  \else
    \draw[darkaccent,line width=0.9pt,fill=white] (\x,0) circle (2pt);
  \fi
  \node[below] at (\x,-0.05) {\lab};
}
\node[above] at (-2.6,0.18) {$+$};
\node[above] at (-1,0.18) {$-$};
\node[above] at (1,0.18) {$+$};
\node[above] at (2.7,0.18) {$-$};
\end{tikzpicture}
\end{center}

\section*{6. Приведение логарифмов к одному виду}

Перед заменой логарифмы должны стать одинаковыми. Полезные формулы:
\[
\log_a uv=\log_a u+\log_a v,
\qquad
\log_a \frac{u}{v}=\log_a u-\log_a v,
\]
\[
\log_{a^m} b=\frac{1}{m}\log_a b,
\qquad
\log_a b=\frac{1}{\log_b a}.
\]

Пример:
\[
\log_{\frac14} A=-\log_4 A.
\]

Если встречается
\[
\log_4(3^x-1)
\]
несколько раз, удобно сделать замену:
\[
t=\log_4(3^x-1).
\]

\section*{7. Чётная степень под логарифмом}

Критически важное правило:
\[
\log_a(x^2)=2\log_a|x|.
\]

Модуль нужен, потому что \(x^2>0\) при \(x\ne0\), но сам \(x\) может быть отрицательным. Без модуля можно получить незаконный логарифм от отрицательного числа.

Более общий вид:
\[
\log_a\bigl(F(x)^{2n}\bigr)=2n\log_a|F(x)|.
\]

\section*{8. Типичные ошибки}

\begin{itemize}
\item Забыть поменять знак при основании \(0<a<1\).
\item Заменить совокупность системой при переходе от \(|F|>G\).
\item Не выколоть нули знаменателя.
\item Потерять ОДЗ логарифма.
\item При вынесении чётной степени из логарифма забыть модуль.
\item Делать замену \(t=\log_a f(x)\), но забыть затем вернуться к \(x\).
\item Писать десятичные приближения там, где лучше оставить дроби или корни.
\end{itemize}

\newpage

\section*{Тренировочный блок}

Решите неравенства. Ответы запишите в виде промежутков.

\begin{enumerate}
\item \(\log_2(x-1)\le3\).
\item \(\log_{\frac13}(x+2)>1\).
\item \(|2x-1|\ge3\).
\item \(\log_2(x^2-5x+6)\ge0\).
\item \(\log_{\frac12}\frac{x-1}{x+2}\le0\).
\item \(\log_3(x-1)+\log_3(x-3)\le1\).
\item \(\log_4(3^x-1)\le\frac12\).
\item \(\log_{\frac14}(2-x)\ge \log_{\frac14}(x+1)\).
\item \(\log_5\bigl((x-1)^2\bigr)\le2\).
\item \(\log_4(3^x-1)\bigl(2-\log_4(3^x-1)\bigr)\le\frac34\).
\end{enumerate}

\newpage

\section*{Ответы}

\renewcommand{\arraystretch}{1.35}
\begin{tabularx}{\linewidth}{@{}cX@{}}
\toprule
№ & Ответ \\
\midrule
1 & \((1;9]\) \\
2 & \((-2;-\frac53)\) \\
3 & \((-\infty;-1]\cup[2;+\infty)\) \\
4 & \(\left(-\infty;\frac{5-\sqrt5}{2}\right]\cup\left[\frac{5+\sqrt5}{2};+\infty\right)\) \\
5 & \((-\infty;-2)\) \\
6 & \((3;4]\) \\
7 & \((0;1]\) \\
8 & \(\left[\frac12;2\right)\) \\
9 & \([-4;1)\cup(1;6]\) \\
10 & \((0;1]\cup[2;+\infty)\) \\
\bottomrule
\end{tabularx}

\end{document}
